\documentclass[12pt,letterpaper]{article}

\usepackage[utf8]{inputenc}
\usepackage[T1]{fontenc}
\usepackage{amsmath}
\usepackage{amssymb}
\usepackage{geometry}
\usepackage[hidelinks]{hyperref}
\usepackage{mathtools}
\usepackage[backend=biber,style=numeric,sorting=none]{biblatex}

\geometry{margin=1in}
\IfFileExists{../../docs/latex/references.bib}{\addbibresource{../../docs/latex/references.bib}}{}

\newcommand{\assoc}{\mathrm{assoc}}
\newcommand{\app}{\mathrm{app}}
\newcommand{\diag}{\mathrm{diag}}
\newcommand{\dd}{\mathrm{d}}
\newcommand{\eff}{\mathrm{eff}}
\newcommand{\mass}{\mathrm{mass}}
\newcommand{\siteD}{\mathrm{D}}
\newcommand{\siteA}{\mathrm{A}}

\title{Explicit Association Closures for PC-SAFT}
\author{}
\date{}

\begin{document}
\maketitle

\section{Purpose and Scope}

This note derives explicit algebraic closures for the association site fractions
used in the PC-SAFT association contribution. The exact PC-SAFT association
model solves nonlinear mass-action equations for the unbonded site fractions.
The closures derived here replace that nonlinear solve by explicit formulas or
by a fixed, unrolled sequence of algebraic maps. Consequently, automatic
differentiation of these closures gives exact derivatives of the approximate
Helmholtz-energy model being evaluated, not exact derivatives of the original
implicit PC-SAFT association model unless the closure is proven to reproduce the
mass-action solution.

The intended first use is a narrow Stage 1 derivation for one-associating-
component Gross--Sadowski 2002 vapor--liquid equilibrium tests, beginning with
methanol/isobutane at \(T=373.15\,\mathrm{K}\) and \(k_{ij}=0.05\)
\cite{Gross2002}. This note does not admit general associating VLE, LLE,
electrolyte, reactive, or temperature-route support.

\section{Exact Association Model}

Let \(i=1,\dots,N_c\) index components and let \(a=1,\dots,N_s\) index the
flattened association sites in the mixture. Site \(a\) belongs to component
\(i(a)\). If a component has \(n_{iA}\) sites of class \(A\), the site-weight
factor for a flattened site \(b\) is
\begin{equation}
  w_b = x_{i(b)}\,\nu_b,
  \label{eq:site_weight}
\end{equation}
where \(x_i\) is the mole fraction and \(\nu_b\) is the multiplicity carried by
flattened site \(b\). In the common convention where each repeated site is
expanded explicitly, \(\nu_b=1\). In a grouped-site convention, \(\nu_b\) carries
the number of equivalent sites represented by \(b\).

For temperature \(T\), molar density \(\rho\), composition vector \(\mathbf{x}\),
and parameter vector \(\mathbf{p}\), let
\(\Delta_{ab}(T,\rho,\mathbf{x},\mathbf{p})\) denote the association strength
between sites \(a\) and \(b\). In PC-SAFT this strength collects the association
energy, association volume, segment diameter, and hard-sphere contact value
terms; its exact implementation follows the PC-SAFT association-strength
relation and the Gross--Sadowski cross-association mixing rules where those
rules are active \cite{Gross2002,Chapman1989,Chapman1990,Huang1990}.

The exact mass-action equation for the unbonded fraction \(X_a\) is
\begin{equation}
  X_a =
  \frac{1}{
    1 + \rho \sum_{b=1}^{N_s} w_b X_b \Delta_{ab}
  }.
  \label{eq:exact_mass_action}
\end{equation}
Equivalently, define
\begin{equation}
  F_a(\mathbf{X};T,\rho,\mathbf{x},\mathbf{p}) =
  X_a\left(1+\rho\sum_{b=1}^{N_s}w_b X_b\Delta_{ab}\right)-1,
  \label{eq:mass_action_residual}
\end{equation}
and solve
\begin{equation}
  \mathbf{F}(\mathbf{X};T,\rho,\mathbf{x},\mathbf{p})=\mathbf{0}.
  \label{eq:mass_action_system}
\end{equation}

The association contribution to the reduced residual Helmholtz energy is
\begin{equation}
  a^\assoc =
  \sum_{i=1}^{N_c} x_i
  \sum_{A\in\Gamma_i} n_{iA}
  \left(
    \ln X_{iA} - \frac{X_{iA}}{2} + \frac{1}{2}
  \right),
  \label{eq:association_helmholtz}
\end{equation}
where \(\Gamma_i\) is the set of site classes on component \(i\). Equation
\eqref{eq:association_helmholtz} is evaluated at the mass-action solution in
the exact model.

\section{Matrix Form}

Define the nonnegative association matrix
\begin{equation}
  K_{ab}(T,\rho,\mathbf{x},\mathbf{p})
  =
  \rho\,w_b\,\Delta_{ab}(T,\rho,\mathbf{x},\mathbf{p}).
  \label{eq:k_matrix}
\end{equation}
Then the exact site balance can be written as
\begin{equation}
  X_a = G_a(\mathbf{X})
  =
  \frac{1}{1+\sum_{b=1}^{N_s}K_{ab}X_b},
  \qquad a=1,\dots,N_s.
  \label{eq:fixed_point_map}
\end{equation}
The exact PC-SAFT association model solves the fixed point of
\eqref{eq:fixed_point_map}. The explicit closures below replace that solve by
closed-form or fixed-depth algebraic approximations
\begin{equation}
  \widetilde{\mathbf{X}}
  =
  \widetilde{\mathbf{X}}(T,\rho,\mathbf{x},\mathbf{p}),
  \label{eq:explicit_x_general}
\end{equation}
and then evaluate
\begin{equation}
  \widetilde{a}^{\assoc} =
  \sum_{i=1}^{N_c} x_i
  \sum_{A\in\Gamma_i} n_{iA}
  \left(
    \ln \widetilde{X}_{iA}
    - \frac{\widetilde{X}_{iA}}{2}
    + \frac{1}{2}
  \right).
  \label{eq:approx_association_helmholtz}
\end{equation}

\section{Closure A: One-Component 2B Sites}
\label{sec:closure-a}

The first exact reduced closure applies to a binary or multicomponent mixture
with only one associating component \(s\), inert nonassociating partner
components, and a 2B donor--acceptor site topology. Let \(n_D\) be the number of
donor-like sites and \(n_A\) the number of acceptor-like sites on component
\(s\). Assume that donor--acceptor association is the only active association
class and that all equivalent donor sites share \(X_D\), all equivalent acceptor
sites share \(X_A\), and the relevant association strength is
\(\Delta_{DA}\). Define
\begin{align}
  c_A &= \rho\,x_s\,n_A\,\Delta_{DA}, \label{eq:c_a}\\
  c_D &= \rho\,x_s\,n_D\,\Delta_{DA}. \label{eq:c_d}
\end{align}
The grouped mass-action equations are
\begin{align}
  X_D &= \frac{1}{1+c_A X_A}, \label{eq:xd_two_class}\\
  X_A &= \frac{1}{1+c_D X_D}. \label{eq:xa_two_class}
\end{align}
Substituting \eqref{eq:xa_two_class} into \eqref{eq:xd_two_class} gives
\begin{align}
  X_D
  &=
  \frac{1}{
    1 + c_A/(1+c_D X_D)
  } \nonumber\\
  &=
  \frac{1+c_D X_D}{1+c_A+c_D X_D}. \label{eq:xd_substitution}
\end{align}
Multiplying through and collecting terms gives the quadratic equation
\begin{equation}
  c_D X_D^2 + (1+c_A-c_D)X_D - 1 = 0.
  \label{eq:xd_quadratic}
\end{equation}
The physical root in \(0<X_D\le 1\) is
\begin{equation}
  X_D =
  \frac{-(1+c_A-c_D)
  + \sqrt{(1+c_A-c_D)^2+4c_D}}
  {2c_D},
  \qquad c_D>0.
  \label{eq:xd_root_direct}
\end{equation}
A numerically preferable algebraically equivalent form is
\begin{equation}
  X_D =
  \frac{2}{
    (1+c_A-c_D)+\sqrt{(1+c_A-c_D)^2+4c_D}
  },
  \label{eq:xd_root_stable}
\end{equation}
with the limiting value \(X_D=(1+c_A)^{-1}\) as \(c_D\to 0\). The acceptor
fraction follows from
\begin{equation}
  X_A = \frac{1}{1+c_D X_D}.
  \label{eq:xa_from_xd}
\end{equation}
For symmetric 2B alcohol-like cases with \(n_D=n_A\), \(c_D=c_A=c\), so
\begin{equation}
  X_D=X_A=\frac{2}{1+\sqrt{1+4c}}.
  \label{eq:symmetric_2b_root}
\end{equation}

For this closure, the approximate Helmholtz expression is exact under the
assumptions listed above:
\begin{equation}
  a^\assoc_{2B}
  =
  x_s\left[
    n_D\left(\ln X_D-\frac{X_D}{2}+\frac{1}{2}\right)
    +
    n_A\left(\ln X_A-\frac{X_A}{2}+\frac{1}{2}\right)
  \right].
  \label{eq:two_class_helmholtz}
\end{equation}
If the implemented site topology or \(\Delta_{ab}\) matrix includes additional
active site-pair interactions, \eqref{eq:two_class_helmholtz} is no longer an
exact reduction of the full mass-action system and must be labeled as an
explicit approximation.

\section{Closure B: Full-Matrix Picard Unroll}
\label{sec:closure-b}

For a general association matrix, define the row-sum strength
\begin{equation}
  S_a = \sum_{b=1}^{N_s} K_{ab}.
  \label{eq:row_sum_strength}
\end{equation}
A bounded algebraic initializer is
\begin{equation}
  X_a^{(0)}
  =
  \frac{2}{1+\sqrt{1+4S_a}}.
  \label{eq:row_sum_initializer}
\end{equation}
This initializer is the symmetric scalar solution for a row-summed association
field and satisfies \(0<X_a^{(0)}\le 1\) for \(S_a\ge 0\).

For a fixed integer \(N_P\), define the Picard proposal
\begin{equation}
  \widehat{X}_a^{(m)}
  =
  \frac{1}{1+\sum_{b=1}^{N_s}K_{ab}X_b^{(m-1)}},
  \qquad m=1,\dots,N_P,
  \label{eq:picard_proposal}
\end{equation}
and either accept it directly,
\begin{equation}
  X_a^{(m)} = \widehat{X}_a^{(m)},
  \label{eq:undamped_picard}
\end{equation}
or use fixed damping,
\begin{equation}
  X_a^{(m)}
  =
  (1-\omega)X_a^{(m-1)}+\omega\widehat{X}_a^{(m)},
  \qquad 0<\omega\le 1.
  \label{eq:damped_picard}
\end{equation}
The explicit closure is the finite composition
\begin{equation}
  \widetilde{X}_a^{\mathrm{Picard}(N_P,\omega)}
  =
  X_a^{(N_P)}.
  \label{eq:picard_closure}
\end{equation}
Recommended diagnostic variants are
\begin{equation}
  (N_P,\omega)\in
  \{(1,1),\ (3,1),\ (3,0.5),\ (5,0.5)\}.
  \label{eq:picard_variants}
\end{equation}
Because \(N_P\) and \(\omega\) are fixed before differentiation,
\eqref{eq:picard_closure} is an explicit algebraic model. Its automatic
derivatives are exact for \(\widetilde{a}^{\assoc}\), not for the implicit
solution unless \(\mathbf{F}(\widetilde{\mathbf{X}})=\mathbf{0}\).

\section{Closure C: Picard Unroll with Diagonal Newton Polish}
\label{sec:closure-c}

Starting from the damped Picard result in \eqref{eq:picard_closure}, compute the
mass-action residual
\begin{equation}
  F_a =
  X_a\left(1+\sum_{b=1}^{N_s}K_{ab}X_b\right)-1.
  \label{eq:polish_residual}
\end{equation}
The exact Jacobian of the residual with respect to the site fractions is
\begin{equation}
  J_{ab}
  =
  \frac{\partial F_a}{\partial X_b}
  =
  \delta_{ab}\left(1+\sum_{c=1}^{N_s}K_{ac}X_c\right)
  +
  X_a K_{ab}.
  \label{eq:mass_action_jacobian}
\end{equation}
A diagonal Newton approximation retains only
\begin{equation}
  D_a =
  J_{aa}
  =
  1+\sum_{c=1}^{N_s}K_{ac}X_c+X_aK_{aa}.
  \label{eq:diagonal_newton_denominator}
\end{equation}
One explicit damped diagonal correction is then
\begin{equation}
  \widetilde{X}_a^{\mathrm{PDN}}
  =
  X_a - \lambda\frac{F_a}{D_a},
  \qquad 0<\lambda\le 1.
  \label{eq:diagonal_newton_update}
\end{equation}
The first candidate for Stage 1 testing is
\begin{equation}
  N_P=3,\qquad \omega=0.5,\qquad \lambda=0.5,
  \label{eq:picard_diag_newton_candidate}
\end{equation}
which may be labeled \texttt{explicit\_picard3\_diag\_newton1}. This
closure preserves the full association matrix, avoids solving a linear system,
and remains a fixed algebraic expression.

\section{Closure D: Collapsed Donor--Acceptor Mean Field}
\label{sec:closure-d}

For inexpensive screening, the site matrix may be collapsed into donor and
acceptor totals. Define
\begin{align}
  C_D &= \sum_{i=1}^{N_c} x_i n_{iD}, \label{eq:total_donor_sites}\\
  C_A &= \sum_{i=1}^{N_c} x_i n_{iA}. \label{eq:total_acceptor_sites}
\end{align}
For \(C_D C_A>0\), define the donor--acceptor mean association strength
\begin{equation}
  \overline{\Delta}_{DA}^{\eff}
  =
  \frac{
    \sum_i\sum_j x_i n_{iD}\,x_j n_{jA}\,\Delta_{iD,jA}
  }{
    C_D C_A
  }.
  \label{eq:mean_field_delta}
\end{equation}
Then set
\begin{align}
  \alpha &= \rho C_A \overline{\Delta}_{DA}^{\eff},
  \label{eq:mean_field_alpha}\\
  \beta &= \rho C_D \overline{\Delta}_{DA}^{\eff}.
  \label{eq:mean_field_beta}
\end{align}
The two-class quadratic becomes
\begin{equation}
  \beta X_D^2 + (1+\alpha-\beta)X_D - 1 = 0,
  \label{eq:mean_field_quadratic}
\end{equation}
with physical root
\begin{equation}
  X_D =
  \frac{2}{
    (1+\alpha-\beta)+\sqrt{(1+\alpha-\beta)^2+4\beta}
  },
  \label{eq:mean_field_xd}
\end{equation}
and
\begin{equation}
  X_A = \frac{1}{1+\beta X_D}.
  \label{eq:mean_field_xa}
\end{equation}
The collapsed Helmholtz contribution is
\begin{equation}
  \widetilde{a}^{\assoc}_{\mathrm{MF}}
  =
  C_D\left(\ln X_D-\frac{X_D}{2}+\frac{1}{2}\right)
  +
  C_A\left(\ln X_A-\frac{X_A}{2}+\frac{1}{2}\right).
  \label{eq:mean_field_helmholtz}
\end{equation}
This closure erases component-specific site fractions and should therefore be
used only as a diagnostic or initializer unless separate validation shows that
the resulting thermodynamic outputs meet the intended tolerance for the target
systems.

\section{Derivative Interpretation}
\label{sec:derivative-interpretation}

Let
\begin{equation}
  \mathbf{q} = (T,\rho,\mathbf{x},\mathbf{p})
  \label{eq:q_vector}
\end{equation}
collect all independent variables and parameters of interest. In the exact
implicit model,
\begin{equation}
  \mathbf{F}(\mathbf{X}^\star(\mathbf{q}),\mathbf{q})=\mathbf{0},
  \label{eq:implicit_solution}
\end{equation}
and the first derivative of the exact site fractions obeys
\begin{equation}
  \frac{\dd\mathbf{X}^\star}{\dd q_k}
  =
  -\mathbf{F}_{\mathbf{X}}^{-1}
  \frac{\partial\mathbf{F}}{\partial q_k}.
  \label{eq:implicit_sensitivity}
\end{equation}
By contrast, the explicit closures define
\begin{equation}
  \widetilde{\mathbf{X}}(\mathbf{q})
  =
  \mathcal{C}(\mathbf{q})
  \label{eq:closure_function}
\end{equation}
for a selected closure \(\mathcal{C}\). The differentiated model is
\begin{equation}
  \widetilde{a}^{\assoc}(\mathbf{q})
  =
  a^{\assoc}(\widetilde{\mathbf{X}}(\mathbf{q}),\mathbf{q}),
  \label{eq:explicit_model}
\end{equation}
and the derivative is the chain rule
\begin{equation}
  \frac{\dd\widetilde{a}^{\assoc}}{\dd q_k}
  =
  \frac{\partial a^{\assoc}}{\partial q_k}
  +
  \sum_{a=1}^{N_s}
  \frac{\partial a^{\assoc}}{\partial X_a}
  \frac{\partial \widetilde{X}_a}{\partial q_k}.
  \label{eq:explicit_chain_rule}
\end{equation}
Automatic differentiation records \(\mathcal{C}\) and therefore gives exact
first and second derivatives of \(\widetilde{a}^{\assoc}\). These derivatives
are not the exact implicit derivatives of \(a^\assoc\) unless
\(\widetilde{\mathbf{X}}=\mathbf{X}^\star\) over the differentiated state.

A useful local error relation follows from a first-order expansion of
\(\mathbf{F}\) around \(\mathbf{X}^\star\):
\begin{equation}
  \widetilde{\mathbf{X}}-\mathbf{X}^\star
  =
  \mathbf{F}_{\mathbf{X}}^{-1}
  \mathbf{F}(\widetilde{\mathbf{X}},\mathbf{q})
  +
  \mathcal{O}\!\left(
    \|\widetilde{\mathbf{X}}-\mathbf{X}^\star\|^2
  \right).
  \label{eq:closure_error_relation}
\end{equation}
Thus, mass-action residual norms, Helmholtz-energy error, fugacity-coefficient
error, pressure error, and equilibrium residual error must be reported together.
Site-fraction error alone is not sufficient to admit an equilibrium route.

\section{Stage 1 Acceptance Diagnostics}

For each explicit closure candidate, report at least
\begin{align}
  e_X &=
  \|\widetilde{\mathbf{X}}-\mathbf{X}^\star\|_{\infty},
  \label{eq:x_error_metric}\\
  r_\mass &=
  \|\mathbf{F}(\widetilde{\mathbf{X}},\mathbf{q})\|_{\infty},
  \label{eq:mass_residual_metric}\\
  e_a &=
  |\widetilde{a}^{\assoc}-a^{\assoc,\star}|,
  \label{eq:a_error_metric}
\end{align}
plus pressure, residual-chemical-potential, fugacity-coefficient, and
equilibrium-residual errors where those quantities are evaluated. For the first
Gross--Sadowski 2002 associating VLE proof, the closure must be exercised first
at EOS states before it is used inside a \(\mathrm{bubble\ pressure}\) solve.

\section{Recommended Stage 1 Order}

\begin{enumerate}
  \item Implement or evaluate Closure A for one-associating-component 2B cases.
  \item Compare Closure A against the exact mass-action solution at pure
        methanol states and methanol/isobutane states from Gross--Sadowski 2002.
  \item Evaluate Closure B and Closure C as explicit approximate models and
        report the diagnostics in Section~\ref{sec:derivative-interpretation}.
  \item Keep Closure D limited to diagnostics and initial values unless output
        errors justify broader use.
  \item Admit no equilibrium route from this derivation alone. The first route
        proof remains the narrow neutral, nonelectrolyte, nonreactive
        \texttt{bubble\_pressure} case with at most one associating
        component.
\end{enumerate}

\printbibliography

\end{document}
